\section*{Hoofdstuk \ref{ch:g8:the-oceans} --- De oceanen}

\begin{solution}{exo:g8:the-oceans:1}
\omterm{def:g8:the-oceans:seafloor}{Continentaal plat}, \omterm{def:g8:the-oceans:seafloor}{continentale helling}, \omterm{def:g8:the-oceans:seafloor}{abyssale vlakte}, \omterm{def:g8:the-oceans:seafloor}{oceanische rug}.
\end{solution}

\begin{solution}{exo:g8:the-oceans:2}
(a) Niet waar: ongeveer \qty{3700}{m}, dat is \qty{3.7}{km}. (b) Waar. (c)
Niet waar: verdamping haalt water weg en laat de zouten achter, dus wordt
het water zouter. (d) Waar.
\end{solution}

\begin{solution}{exo:g8:the-oceans:3}
(a) Een \omterm{def:g8:the-oceans:seafloor}{oceanische rug}. (b) Het \omterm{def:g8:the-oceans:salinity}{zoutgehalte}. (c) Een \omterm{def:g8:the-oceans:current}{gyre}. (d) Een \omterm{def:g8:the-oceans:seafloor}{diepzeetrog}.
\end{solution}

\begin{solution}{exo:g8:the-oceans:4}
(a) Tegen de klok in. (b) Warm: de Golfstroom; koud: de Canarische Stroom.
(c) De \omterm{def:g8:winds-global-circulation:global-winds}{passaatwinden}.
\end{solution}

\begin{solution}{exo:g8:the-oceans:5}
(a) $35 \div 1000 \times 100 = 3.5~\%$. (b) $20 \times 35 = \qty{700}{g}$.
(c) $1000 \div 35 \approx 28.6$: ongeveer \qty{29}{kg}.
\end{solution}

\begin{solution}{exo:g8:the-oceans:6}
(a) $1500 \times 6 = \qty{9000}{m}$. (b) Het geluid gaat naar de bodem en
komt weer omhoog: het legt de diepte twee keer af. (c) $9000 \div 2 =
\qty{4500}{m}$. (d) Een \omterm{def:g8:the-oceans:seafloor}{abyssale vlakte} (een trog zou dieper zijn, een rug
of een plat minder diep).
\end{solution}

\begin{solution}{exo:g8:the-oceans:7}
(a) Tussen ongeveer \qty{3500}{m} en \qty{5000}{m}. (b) Ongeveer
\qty{4400}{km} langs het profiel, ongeveer \qty{2300}{m} diep. (c) Ongeveer
$4500 - 2300 = \qty{2200}{m}$ (tussen 1200 en \qty{2700}{m}, afhankelijk
van de gekozen \omterm{def:g1:landscapes-around-us:landform}{vlakte}). (d) De \omterm{def:g2:water-on-earth:ocean}{oceaan} is een paar kilometer diep maar
duizenden kilometers breed: op dezelfde schaal getekend zou het hele
profiel dunner zijn dan een lijn.
\end{solution}

\begin{solution}{exo:g8:the-oceans:8}
\omterm{def:g2:water-on-earth:ocean}{Zee} A is zouter: de zon en de droge lucht veroorzaken sterke
\emph{verdamping}, het water verdwijnt, de \emph{zouten blijven achter}, en
bijna geen \omterm{def:g4:clouds-rain-snow:precipitation}{regen} of \omterm{def:g1:landscapes-around-us:water}{rivier} brengt \emph{\omterm{def:g2:water-on-earth:salt-fresh}{zoet water}} terug. \omterm{def:g2:water-on-earth:ocean}{Zee} B is minder
zout: haar \omterm{def:g1:landscapes-around-us:water}{rivieren} en haar \omterm{def:g4:clouds-rain-snow:precipitation}{regen} brengen veel \emph{\omterm{def:g2:water-on-earth:salt-fresh}{zoet water}}, dat de
zouten verdunt.
\end{solution}

\begin{solution}{exo:g8:the-oceans:9}
(a) De Golfstroom, eerst naar het noorden langs de \omterm{def:g1:landscapes-around-us:water}{kust}, dan naar het
noordoosten. (b) Naar het oosten (noordoosten), met de Noord-Atlantische
Stroom. (c) Langs de \omterm{def:g1:landscapes-around-us:water}{kust} van Noordwest-Afrika (en van Spanje en
Portugal), met de Canarische Stroom. (d) Koud.
\end{solution}

\begin{solution}{exo:g8:the-oceans:10}
(a) $6.4 \times 24 = 153.6$: ongeveer \qty{154}{km}. (b) $153.6 \times 7
\approx \qty{1075}{km}$. (c) De Golfstroom is iets sneller dan een
wandelaar, maar hij stopt nooit, dag en nacht.
\end{solution}

\begin{solution}{exo:g8:the-oceans:11}
(a) \qty{3.7}{km} en \qty{11}{km}. (b) $11000 \div 3700 \approx 3.0$
keer. (c) Het geluid legt $2 \times 11000 = \qty{22000}{m}$ af; $22000
\div 1500 \approx 14.7$: ongeveer \qty{15}{s}. (d) $1000 = 10^3$ jaar.
\end{solution}

\begin{solution}{exo:g8:the-oceans:12}
(a) Bij de \omterm{def:g3:continents-and-oceans:equator}{evenaar} duwen de \omterm{def:g8:winds-global-circulation:global-winds}{passaatwinden} warm tropisch water naar het
westen. Aan de westrand van de Atlantische \omterm{def:g2:water-on-earth:ocean}{Oceaan} buigt het naar het
noorden af met de Golfstroom, een warme stroming, langs Florida en
Noord-Amerika. De \omterm{def:g8:winds-global-circulation:global-winds}{westenwinden} duwen het dan over de Atlantische \omterm{def:g2:water-on-earth:ocean}{Oceaan}
naar het noordoosten, tot in de \omterm{def:g2:water-on-earth:ocean}{zeeën} voor Noorwegen. (b) Daar geeft het
zijn warmte af aan de koude lucht: het koelt af, en het wordt zouter als er
zee-ijs ontstaat, omdat het zout in het water achterblijft. Koud, \omterm{def:g2:water-on-earth:salt-fresh}{zout
water} is dicht, dus zakt het naar de bodem. (c) Het kruipt langs de
oceaanbodem naar het zuiden, naar Antarctica, en dan de Indische en de
Stille \omterm{def:g2:water-on-earth:ocean}{Oceaan} in, waar het langzaam weer opstijgt.
\end{solution}

\begin{solution}{pb:g8:the-oceans:1}
\textbf{1.} In de Atlantische \omterm{def:g2:water-on-earth:ocean}{Oceaan}, op het \omterm{def:g3:continents-and-oceans:equator}{noordelijk halfrond}.

\textbf{2.} Met de klok mee; de boei wordt eerst naar het noorden meegevoerd.

\textbf{3.} Etappe A (Golfstroom) is warm: hij komt uit de tropen. Etappe
C (Canarische Stroom) is koud: hij komt van hogere breedten, op weg naar de
\omterm{def:g3:continents-and-oceans:equator}{evenaar}.

\textbf{4.} Etappe B: de \omterm{def:g8:winds-global-circulation:global-winds}{westenwinden}. Etappe D: de \omterm{def:g8:winds-global-circulation:global-winds}{passaatwinden}.

\textbf{5.} $3800 \div 100 = 38$ dagen.

\textbf{6.} B: $2100 \div 20 = 105$ dagen. C: $3100 \div 20 = 155$ dagen.
D: $6100 \div 20 = 305$ dagen.

\textbf{7.} $3800 + 2100 + 3100 + 6100 = \qty{15100}{km}$.

\textbf{8.} De Golfstroom is de smalle, snelle stroming aan de westkant van
de \omterm{def:g8:the-oceans:current}{gyre}; de stromingen aan de andere kanten zijn breed en langzaam.

\textbf{9.} $38 + 105 + 155 + 305 = 603$ dagen.

\textbf{10.} $603 \div 30 = 20.1$: ongeveer \textbf{20 maanden}.

\textbf{11.} Na de etappes A, B en C zijn er $38 + 105 + 155 = 298$ dagen
verstreken; de 365 dagen vallen tijdens etappe D (de Noordequatoriale
Stroom), $365 -
298 = 67$ dagen na het begin ervan.

\textbf{12.} $603 - 365 = 238$ dagen meer, dat is $238 \div 30 \approx 8$
maanden meer.
\end{solution}
